% !TeX root = ../本科微分几何笔记.tex
% 由主文档通过 \input 载入；请编译主文档。
\setcounter{chapter}{1} % 保留原书章号。
\chapter{符号、公式与学习检查}
\label{app:dg-reference}

\section{主要符号与适用条件}
\begin{longtable}{p{0.16\textwidth}p{0.36\textwidth}p{0.38\textwidth}}
\toprule
符号 & 含义 & 使用时应检查的条件\\
\midrule
\endhead
$\gamma(t),s$ & 参数曲线、弧长参数 & 弧长换元要求 $\gamma'\ne0$\\
$\kappa,\tau$ & 无符号曲率、挠率 & Frenet 挠率要求 $\kappa>0$\\
$k,\theta$ & 平面有符号曲率、切向角 & 平面定向与行进方向固定\\
$\mathbf x(u,v),\Sigma$ & 参数片、曲面 & 秩为 $2$，参数片另需同胚条件\\
$T_p\Sigma,\nu$ & 切空间、单位法向 & 法向可只局部存在\\
$E,F,G$ & 第一基本形式系数 & $E>0,\ EG-F^2>0$\\
$e,f,g$ & 第二基本形式系数 & 依赖法向选择\\
$g_{ij},g^{ij}$ & 度量矩阵及其逆的系数 & 上指标不是逐项倒数\\
$S,k_1,k_2$ & 形算子、主曲率 & $S=-\dd\nu$，特征值相对于度量\\
$K,H$ & 高斯曲率、平均曲率 & 反向法向使 $H$ 变号，$K$ 不变\\
$\Gamma^k_{ij},\nabla$ & Christoffel 符号、协变导数 & Christoffel 符号本身不是张量\\
$k_g,k_n,\tau_g$ & 测地曲率、法曲率、测地挠率 & 使用单位速 Darboux 标架\\
$\omega,\dd A$ & 联络一形式、面积元或形式 & $\dd\omega=-K\dd A$ 中使用定向形式\\
$\chi,g$ & 欧拉示性数、亏格 & $\chi=2-2g$ 用于连通闭可定向曲面\\
\bottomrule
\end{longtable}
符号 $g$ 在不同语境中可表示第二基本形式系数、抽象度量或亏格；每次出现均以所在章节的定义为准。同一推导中用 $g_{ij}$ 或 $\mathbf g$ 表示度量矩阵，避免与第二基本形式的末项系数混淆。

\section{常用公式的查阅路径}
曲线的一般参数曲率与挠率见式~\eqref{eq:dg-curvature-general}、\eqref{eq:dg-torsion-general}，Frenet 公式见式~\eqref{eq:dg-frenet}。曲面长度与面积分别见式~\eqref{eq:dg-surface-length}、\eqref{eq:dg-area}；形算子、$K,H$ 与 Euler 法曲率公式分别见式~\eqref{eq:dg-shape-matrix}、\eqref{eq:dg-KH}、\eqref{eq:dg-euler-curvature}。图像曲面的计算可从式~\eqref{eq:dg-graph-K}、\eqref{eq:dg-graph-H} 开始，但必须先选定法向。

内蕴计算先由式~\eqref{eq:dg-christoffel} 求 Christoffel 符号，再根据问题选择式~\eqref{eq:dg-geodesic-equations} 的测地线方程或式~\eqref{eq:dg-K-christoffel} 的曲率公式。测地曲率与角度的关系见式~\eqref{eq:dg-kg-angle}，局部、整体及带边界的 Gauss--Bonnet 公式分别见式~\eqref{eq:dg-local-gb}、\eqref{eq:dg-global-gb}、\eqref{eq:dg-boundary-gb}。公式使用的区域、正则性和定向条件以正文命题为准。

\section{综合计算的一般顺序}
给定一个参数曲面后，先明确参数域并检查 $\mathbf x_u\times\mathbf x_v\ne0$；若讨论整个几何对象，再检查坐标是否遗漏点、是否重复覆盖。随后计算 $E,F,G$，确认 $EG-F^2>0$，再指定单位法向并计算 $e,f,g$。主曲率应由 $\det(\mathbf b-k\mathbf g)=0$ 求出，并用和为 $2H$、积为 $K$ 相互核对。曲率符号可用平面、球面、圆柱及鞍面四类模型检验。

若研究内蕴问题，应优先利用度量中的对称性：常系数度量对应平直坐标，旋转曲面可以使用 Clairaut 关系，保角度量可直接使用二维保角曲率公式。声称一个映射等距时，应验证整个切向量内积，或等价地验证拉回后的度量矩阵，而不是仅检查某条曲线的长度。声称最短时，还需区分测地线方程、局部最短性和整体最短性。

\section{证明依赖与课程完成标准}
本书直接证明的主线包括弧长与 Frenet 理论、曲线论基本定理、基本形式的坐标不变性、形算子自伴性、曲面基本方程、Gauss 绝妙定理、面积和能量第一变分、测地线局部理论以及在明确拓扑前置结果下的 Gauss--Bonnet 推导。逆函数与隐函数定理、常微分方程初值定理、谱定理、Green--Stokes 公式属于先修内容；等周不等式的证明另外引用 Parseval 恒等式。平面转角定理、曲面三角剖分与分类定理作为拓扑结论引用，曲面论基本定理的存在性部分给出依赖可积性定理的证明思路，Hopf--Rinow 定理仅陈述，不伪装为已在本书证明。

完成本书核心学习后，应能从具体参数化出发计算曲率、识别方向与坐标陷阱，解释第一基本形式与第二基本形式各自记录的信息，独立推导常见曲面的测地线方程，并在满足条件的曲面区域上使用 Gauss--Bonnet 公式。进一步的目标是把这些计算重述为与坐标无关的线性映射、双线性型和联络结论，从而自然进入流形与黎曼几何。

